\documentclass[../main.tex]{subfiles}

\begin{document}
\subsection{{\Statement} {\Control}}
\subsubsection{Percabangan \textit{If}}

\paragraph{\textit{If} Tunggal: mengerjakan perintah pada cabang program hanya jika kondisi terpenuhi.}

\begin{code}[r]
\begin{verbatim}
if (kondisi) { statemen }
\end{verbatim}
\end{code}

\paragraph{\textit{If-else}: mengerjakan perintah pada kondisi \textit{if} bernilai benar atau \textit{if} bernilai salah yang ditetapkan sebagai \textit{else}.}

\begin{code}[r]
\begin{verbatim}
if (kondisi) { statemen }
else (kondisi) { statemen }
\end{verbatim}
\end{code}

\paragraph{\textit{Nested if}: percabangan pada \textit{if-else} yang memiliki lebih dari 2 pilihan.}

\begin{code}[r]
\begin{verbatim}
if (kondisi) { statemen }
else if (kondisi) { statemen }
else (kondisi) { statemen }
\end{verbatim}
\end{code}

\subsubsection{\textit{Switch case}: merupakan bentuk lain dari \textit{nested if}, hal ini karena \textit{switch case} memiliki prinsip kerja yang sama, tetapi dengan format penulisan yang berbeda \parencite{hanief2020konsep}.}

\begin{code}[n]
\begin{verbatim}
switch (variabel) {
   case 0:
       statemen
       break;
   case 1:
       statemen
       break;
   ...
   case n:
       statemen
       break;
   default:
       statemen
       break;
}
\end{verbatim}
\end{code}

\subsection{{\Looping}}
Struktur perulangan dapat dibedakan menjadi tiga jenis yaitu \textit{for}, \textit{while} dan \textit{dowhile} \parencite{anam2021cara}.

\stepcounter{subsubsection}
\begin{enumerate}[labelsep=2.8ex]
    \item {\textit{For}: perulangan yang batasan dan syarat pengulangannya sudah ditentukan \parencite{suryadi2018belajar}.}
\end{enumerate}

\begin{code}[n]
\begin{verbatim}
 for (inisialisasi; kondisi perulangan; pengubah nilai cacah)
\end{verbatim}
\end{code}

\subsubsection{\textit{While}: perulangan yang akan dijalankan selama kondisi terpenuhi dengan pemeriksaan kondisi sebelum dijalankan \parencite{anam2021cara}.}

\begin{code}[n]
\begin{verbatim}
while (kondisi) {
   statemen
}
\end{verbatim}
\end{code}

\subsubsection{\textit{Do-while}: perulangan yang dijalankan terlebih dahulu 1 kali, kemudian akan ada pemeriksaan kondisi \textit{while}, jika terpenuhi perulangan terjadi \parencite{anam2021cara}.}

\begin{code}[n]
\begin{verbatim}
do {
    statemen
} while (kondisi)
\end{verbatim}
\end{code}
\end{document}
